\section{EPGSpeller Recognition System}
\label{sec:system}

This section defines the EPGSpeller recognition system.
EPGSpeller provides a recognition layer for greedy character-string decoding on top of the input modality demonstrated by SilentSpeller~\cite{kimura2022silentspeller}.

\subsection{Overview and CTC Decoding}

EPGSpeller is a sequence recognizer that takes electropalatographic contact frame sequences and outputs the corresponding character strings.
Each trial is a variable-length contact sequence; character boundaries and phoneme boundaries are not provided.
The primary output path is lexicon-free greedy character-string decoding defined in Section~\ref{sec:task:definition}.

In silent-spelling EPG sequences, the correspondence between each frame and a character is not explicit.
Participants articulate continuously, and tongue--palate contact is not segmented at the character level.
EPGSpeller uses connectionist temporal classification (CTC) for this alignment~\cite{graves2006connectionist}.
CTC learns sequence outputs without frame-level labels by collapsing frame-level outputs that include blank symbols and repetitions into the final character string.

\begin{table}[!b]
\centering
\caption{Boundary between main recognizer and design analyses.}
\label{tab:recognizer_design}
\setlength{\tabcolsep}{1.5pt}
\renewcommand{\arraystretch}{1.03}
\begin{tabular}{@{}p{0.34\columnwidth}p{0.18\columnwidth}p{0.38\columnwidth}@{}}
\toprule
Element & Claims? & Role \\
\midrule
vector/GRU + greedy decode & Yes & Primary result \\
LEX projection & No & Closed-vocab.\ diag. \\
Spatial / electrode variants & No & Engineering analyses \\
Preprocessing controls & No & Transfer-gap control \\
\bottomrule
\end{tabular}
\end{table}

\subsection{Main Recognizer: Vector/GRU Path}

In the main recognizer, each electropalatographic frame is treated as a flat contact vector.
The temporal recognizer is a GRU-based recurrent CTC decoder.
We treat this vector/GRU CTC path as the main recognizer throughout the paper.
The central claim that electropalatographic silent spelling can support greedy character-string decoding when personalized calibration is available is evaluated on this main recognizer and greedy decoding.

Multichannel contact frames are mapped through an input projection into a representation that the sequence recognizer can process efficiently.
Along the time axis, compression reduces sequence length while preserving local context.
During training, we apply dropout, noise injection, smoothing, and SpecAugment-style temporal and channel regularization to improve robustness to contact-position variability, missing channels, and participant-specific differences~\cite{park2019specaugment}.

\subsection{Spatial Front Ends and Electrode Reduction}

SmartPalate contact channels have a layout on the palate.
We compare row-column, grid, and patch front ends to test whether spatial structure helps electropalatographic silent-spelling recognition.
The row-column front end summarizes contact patterns on the surrogate palatal grid into row-wise and column-wise activity.
Grid and patch front ends use local spatial structure more explicitly, but electropalatographic contact is sparse and strongly participant-dependent, so a strong two-dimensional spatial assumption is not always advantageous.
Results in Section~\ref{sec:results} show that the flat vector/GRU path remains competitive on P1.

We also compare contact-frequency-based selection, spatially dispersed selection, ranking transfer from other participants, and random subsets to examine whether recognition performance can be preserved with fewer channels.
Electrode reduction and missing-channel robustness are engineering analyses; they are not primary means of solving cross-user transfer.
Later analyses indicate that K=64 channels retain near-main P1/P2 performance for the selection methods tested.

\subsection{Implementation Details}

The main recognizer connects a CTC output layer to a bidirectional GRU (1024 units, 5 layers); input projection maps contact-channel representations to 256 dimensions.
Along time, compression with kernel length~32 and stride~4 reduces sequence length.
The output has 40 classes (26 alphabet letters, blank, CTC blank, and other special symbols).
Training uses batch size~64, fixed learning rate~0.02, L2 regularization coefficient 1e-5, dropout rate~0.4, and 10{,}000 training steps.
Data augmentation applies SpecAugment-style temporal masking (T=12, rate~0.4), electrode masking (F=6, rate~0.4), and time warping (W=0.05, rate~0.2), plus noise and smoothing.
Results are averaged over four trials with random seeds 0--3; mean and standard deviation are reported.
In PCA variants, 120 contact channels are compressed to 16 principal components before input projection.
Key hyperparameters are listed in the supplementary configuration table.

Spatial front ends, electrode reduction, and regularization are engineering design comparisons for representing and simplifying SmartPalate-style electropalatographic signals.
Table~\ref{tab:recognizer_design} separates the main recognizer from auxiliary analyses referenced later in the paper.
The next section defines the evaluation conditions under which this recognizer is compared.
